{
\footnotesize
\setlength{\tabcolsep}{3pt}
\Needspace{13\baselineskip}
\begin{longtable}{|>{\centering\arraybackslash}m{1cm}|m{2.5cm}|m{5.1cm}|m{3cm}|>{\centering\arraybackslash}m{1.2cm}|}
\caption{Kebutuhan Nonfungsional Sistem}
\label{tab:kebutuhan_nonfungsional} \\
\hline
\multicolumn{1}{|>{\centering\arraybackslash}m{1cm}|}{\textbf{ID}} & \multicolumn{1}{>{\centering\arraybackslash}m{2.5cm}|}{\textbf{Kebutuhan Nonfungsional}} & \multicolumn{1}{>{\centering\arraybackslash}m{5.1cm}|}{\textbf{Deskripsi}} & \multicolumn{1}{>{\centering\arraybackslash}m{3cm}|}{\textbf{Target Metrik}} & \multicolumn{1}{>{\centering\arraybackslash}m{1.2cm}|}{\textbf{Tujuan}} \\
\hline
\endfirsthead
\hline
\multicolumn{1}{|>{\centering\arraybackslash}m{1cm}|}{\textbf{ID}} & \multicolumn{1}{>{\centering\arraybackslash}m{2.5cm}|}{\textbf{Kebutuhan Nonfungsional}} & \multicolumn{1}{>{\centering\arraybackslash}m{5.1cm}|}{\textbf{Deskripsi}} & \multicolumn{1}{>{\centering\arraybackslash}m{3cm}|}{\textbf{Target Metrik}} & \multicolumn{1}{>{\centering\arraybackslash}m{1.2cm}|}{\textbf{Tujuan}} \\
\hline
\endhead
\endfoot
\hline
\endlastfoot

KNF-01 & Kinerja Latensi Rendah & \textit{Online feature serving} berbasis Valkey dan \textit{vector similarity search} berbasis Qdrant harus memberikan latensi yang stabil dan rendah untuk mendukung \textit{inference} waktu nyata pada berbagai domain aplikasi, termasuk pada saat beban tinggi. & Latensi p99 sub-detik indikatif pada lingkungan \textit{node} tunggal untuk \textit{feature serving} dan \textit{vector search}, dengan ambang spesifik per layanan diukur dan dipantau melalui SLO Sloth & T-4 \\
\hline
KNF-02 & \textit{Throughput} Tinggi & Sistem harus mampu menangani \textit{workload} dengan \textit{throughput} tinggi yang umum pada beban kerja \textit{streaming} dan \textit{serving} bervolume tinggi, dengan karakteristik skalabilitas mendekati linear ketika sumber daya ditambah. & \textit{Throughput} \textit{ingest} dan \textit{serving} tumbuh proporsional terhadap jumlah \textit{partition} Kafka serta \textit{replica} Feast melalui \textit{autoscaling} KEDA dan HPA & T-4 \\
\hline
KNF-03 & Skalabilitas Horizontal & Seluruh komponen sistem harus mendukung skala horizontal, sehingga kapasitas dapat ditingkatkan dengan menambah \textit{node} dan menghasilkan peningkatan kinerja yang linear atau mendekati linear. & Seluruh \textit{workload} terstandardisasi pada HPA (CPU/memori), VPA \textit{recommender}, dan KEDA (Kafka \textit{lag}, \textit{custom metrics}) & T-1 \\
\hline
KNF-04 & Keandalan \& Toleransi Kesalahan & Sistem harus andal dengan komponen \textit{stateful} (Kafka, ClickHouse, PostgreSQL, MinIO, Valkey, Qdrant) yang dapat berjalan dalam mode \textit{multi-replica} pada lingkungan produksi multi-\textit{node}, dan mampu melakukan pemulihan otomatis dari kegagalan sementara melalui mekanisme \textit{operator} dan \textit{snapshot} berkala. Pada lingkungan verifikasi \textit{node} tunggal, \textit{replica} \textit{idle} ditetapkan menjadi satu dan ditingkatkan secara dinamis ketika beban nyata terdeteksi. & SLO \textit{platform} dijaga melalui Sloth (Kafka 99,9\%, Postgres 99,95\%, MinIO p99 <500ms pada 99\%), sedangkan \textit{disaster recovery} berbasis Velero \textit{snapshot} harian penuh dan \textit{incremental} per \textit{use case} & T-1 \\
\hline
KNF-05 & Konsistensi \& Kebenaran & Sistem harus menegakkan kontrak kualitas data sepanjang \textit{pipeline}, mencakup kesesuaian skema, validitas rentang nilai, dan kelengkapan data, sehingga hasil analisis maupun \textit{inference} model bertumpu pada data yang sahih dan konsisten. & Seluruh \textit{checkpoint} Great Expectations lolos tanpa pelanggaran skema maupun rentang nilai pada data yang dipromosikan ke \textit{gold layer} & T-2 \\
\hline
KNF-06 & \textit{Observability} & Sistem harus memiliki \textit{observability} menyeluruh, mencakup metrik (Prometheus dan Grafana), \textit{structured log} (Loki), \textit{distributed tracing} (Grafana Tempo melalui OpenTelemetry), \textit{continuous profiling} (Pyroscope), dan \textit{alerting} (Alertmanager dan Sloth) sehingga perilaku sistem dapat dianalisis dari ujung ke ujung. & Seluruh komponen terinstrumentasi metrik dan log standar, \textit{trace} OpenTelemetry untuk lintasan kritis, serta SLO Sloth untuk indikator layanan & T-2 \\
\hline
KNF-07 & Keamanan & Sistem harus menerapkan kontrol keamanan menyeluruh, mencakup autentikasi tunggal melalui dex sebagai \textit{identity provider} OIDC dan oauth2-proxy, otorisasi berbasis RBAC dan SpiceDB, enkripsi data \textit{in transit} dengan Istio mTLS \textit{strict}, enkripsi \textit{at rest} pada MinIO melalui KES sebagai \textit{KMS gateway} ke OpenBao, kebijakan admisi Kyverno, serta deteksi ancaman \textit{runtime} oleh Falco dan pemindaian kerentanan \textit{image} oleh Trivy Operator. & TLS 1.3 \textit{in transit}, AES-256 \textit{at rest}, dan \textit{audit log} OpenBao serta Kubernetes audit \textit{policy} aktif pada seluruh \textit{control plane} & T-2 \\
\hline
KNF-08 & Ekstensibilitas \& Modularitas & Arsitektur harus modular sehingga komponen dapat diperluas atau diganti melalui antarmuka terdefinisi (\textit{Custom Resource Definition}, helm \textit{chart}, dan \textit{overlay} Kustomize) tanpa mengganggu komponen lain. & Setiap komponen disusun pada \textit{folder} terpisah dengan \textit{kustomization.yaml} sendiri, sehingga dapat dipasang, diperbarui, atau diganti secara independen & T-1 \\
\hline
KNF-09 & Akses Konsol Terpadu & Antarmuka konsol pengembang dan operator platform harus dapat diakses dari satu titik autentikasi tunggal sehingga konsol Argo CD, Argo Workflows, Kubeflow, MLflow, DataHub, dan Grafana terhubung melalui satu sesi identitas, dengan dokumentasi referensi yang ringkas untuk konfigurasi setiap komponen. & Seluruh antarmuka konsol platform terhubung pada satu jalur autentikasi dex dan oauth2-proxy tanpa konfigurasi identitas tambahan per layanan, serta tiap komponen memiliki dokumentasi referensi pada repositori Gitea & T-1 \\
\hline
KNF-10 & Efisiensi Sumber Daya & Sistem harus memanfaatkan sumber daya komputasi dan penyimpanan secara efisien untuk mengendalikan biaya tanpa mengorbankan kinerja, baik pada lingkungan \textit{node} tunggal maupun multi-\textit{node}. & Kompresi kolumnar pada ClickHouse, \textit{autoscaling} berbasis HPA, VPA, dan KEDA agar utilisasi sumber daya proporsional terhadap beban, serta atribusi biaya per \textit{workload} melalui OpenCost & T-1 \\
\hline
KNF-11 & Portabilitas & Sistem harus dapat dijalankan di berbagai lingkungan \textit{deployment} dengan perubahan terbatas pada konfigurasi dan tanpa modifikasi kode besar, sebab seluruh komponen berdiri di atas distribusi Kubernetes standar dan \textit{image} \textit{open source}. & Dapat di-\textit{deploy} pada distribusi Kubernetes \textit{managed} (AWS EKS, GCP GKE, Azure AKS) maupun \textit{on-premise} (k3s, kubeadm) hanya dengan penyesuaian \textit{overlay} Kustomize dan parameter \textit{secret} & T-1 \\
\hline
KNF-12 & Kemudahan Pemeliharaan & Sistem harus mudah dipelihara, dengan struktur kode dan manifes yang jelas, \textit{pipeline} CI yang otomatis pada Tekton, serta \textit{deployment} deklaratif dari sumber kebenaran tunggal pada Gitea melalui Argo CD. & Setiap komponen dapat direkonsiliasi ulang dari \textit{repository} Gitea tanpa intervensi manual, dan perubahan ditolak otomatis oleh \textit{pull request review} dan \textit{policy} Kyverno bila tidak sesuai konvensi & T-1 \\
\hline
\end{longtable}
}